\color{unimain}{\subsection{Токовая защита нулевой последовательности компенсационной обмотки}}
\color{black}

\begin{enumerate}

\item
Функция токовой защиты нулевой последовательности компенсационной обмотки реактора, предназначенных для защиты от витковых замыканий в компенсационной обмотки, а также для резервирования действия основных быстродействующих защит при внутренних коротких замыканий. 
\item
Защита выполнена двухступенчатой и реагирует расчетное значение утроенного тока нулевой последовательности, определяемого по формуле:
\begin{equation}
3\dot{I}_{0} = \dot{I}_{A} + a^2 \cdot \dot{I}_{B} + a \cdot \dot{I}_{C}
\end{equation}

\item
Каждая из ступеней имеет регулируемые значения по току и времени срабатывания. Каждая из ступеней защиты может выполнена с блокировкой по второй гармонике (отношению уровня второй гармоники к первой). Данная блокировка реализована на базе модуля \textbf{«ОВ БНТ»}. Предусмотрена блокировка по току 3I0 стороны ВН реактора.

\item
Алгоритм работы ТЗНП КО-1 приведен на рисунке \ref{pict_tznpko:TZNP_KO}, алгоритм работы ТЗНП КО-2 выполнен аналогично. Уставки ТЗНП КО приведены в таблице \ref{app_set:tznpko}.

\medskip
\begin{figure}[H]
\centering
\includegraphics[width=0.9\textwidth,height=0.9\textheight,keepaspectratio]{\fbpath/11.02 УШР/ТЗНП КО/img/TZNP_KO.pdf}
\captionof{figure}{Функциональная схема логики ТЗНП КО-1}\label{pict_tznpko:TZNP_KO}
\end{figure}


\end{enumerate}